% Section "Setting and Notation" (owner: prelim). Definitions plus three proved statements; no empirical result appears here.
% Labels provided: sec:prelim, eq:logit, eq:decomp, lem:invariance (= prop:monotone, kept as an alias), prop:itemdep,
% prop:contam. Numbering follows acmart (theorem counter within the section).
% The binormal statement (Proposition 3) is self-contained: delete everything between PROP3-BEGIN and PROP3-END if space is
% needed (about a third of a page); nothing else in this file depends on it.
% No macros are defined here; the section uses only the packages loaded by main.tex (amsmath, acmart's newtxmath).
\section{Setting and Notation}
\label{sec:prelim}

\paragraph{Yes/no recommender and token confidence.}
A pointwise yes/no recommender~\citep{bao2023tallrec,zhang2025collm} scores a user $u$ (history $H_u$) and a candidate
item $i$ by asking a binary question $q$ about the pair (default $q=\textsf{like}$, ``Would this user like the candidate
item?'') and reading the next-token distribution of a language model $f$, zero-shot or with a LoRA adapter fitted to
yes/no labels. Its \emph{confidence} is the log-odds
\begin{equation}\label{eq:logit}
  \ell(u,i)=\log P_f(\texttt{Yes}\mid x_q(u,i))-\log P_f(\texttt{No}\mid x_q(u,i)),
\end{equation}
with $\hat p=\sigma(\ell)$, where $P_f(\texttt{Yes}\mid x)$ sums the probabilities of the vocabulary ids whose decoded text,
stripped and lower-cased, is ``yes'' (likewise \texttt{No}): the answer is read from token ids, not parsed from generated
text, and candidates are ranked by $\ell$.

\paragraph{Panels and metrics.}
A \emph{unit} is a user (rated panels) or a prediction event (next-item panels) with candidate set $C_u$ and labels
$y(u,i)\in\{0,1\}$. \emph{Rated panels} take the user's earlier rated items, shown with their ratings, as history and the
user's own later rated items as candidates, with $y=1$ for rating $\ge4$, $y=0$ for rating $\le2$ and 3-star items dropped,
so the label answers the question asked. With $c(x)=1,\tfrac12,0$ for $x>0,\,x=0,\,x<0$, the per-user AUC of a score $s$ and
its mean over users with both classes (UAUC~\citep{zhang2025collm,zhang2024binllm}) are
\[
  \mathrm{AUC}_u(s)=\frac{1}{|C_u^+||C_u^-|}\sum_{i\in C_u^+}\sum_{j\in C_u^-}c(s_i-s_j),\quad
  \mathrm{UAUC}=\operatorname*{mean}_{u}\mathrm{AUC}_u .
\]
\emph{Same-candidate next-item panels} hold, per event, one positive and 100 negatives sampled by popularity, ranked
identically by every system (this is not full-catalogue ranking).% TODO: cite Krichene and Rendle (sampled metrics) once in the bib
\ If $a$ candidates score above the positive and $t$ tie with it, its tie-aware rank $R$ is uniform on
$\{1{+}a,\dots,1{+}a{+}t\}$; HR@$k$ is the expectation of $\mathbf 1[R{\le}k]$, NDCG@$k$ of $\mathbf 1[R{\le}k]/\log_2(1{+}R)$
and MRR of $1/R$ (we report $k{=}10$). Design decisions use \emph{dev} users only; every confirmatory estimate uses fresh
\emph{confirmation} users, disjoint by id.

\paragraph{Decomposition and audit quantities.}
Let $\pi(i)=\mathbb E_D[\ell(D,i)]$ be the \emph{item prior}, the expected confidence for $i$ under the history of a donor
$D$ who has not seen $i$ (estimated by the mean $\hat\pi(i)$ over $K$ donors). With the user offset
$b_u=|C_u|^{-1}\sum_{i\in C_u}[\ell(u,i)-\pi(i)]$, the remainder is the \emph{personal evidence} $e$:
\begin{equation}\label{eq:decomp}
  \ell(u,i)=b_u+\pi(i)+e(u,i),\qquad \textstyle\sum_{i\in C_u}e(u,i)=0,
\end{equation}
so within a unit only $\pi+e$ orders candidates. The audit reports three quantities. The \emph{personal share}
$\mathcal S=1-\mathrm{cor}(\tilde\ell,\tilde\pi)^2$, with $\mathrm{cor}$ pooled over pairs after centring within units, is
the fraction of within-unit variance of $\ell$ that no linear function of the prior explains; as $\hat\pi$ is noisy,
$\mathrm{cor}^2$ is divided by the Spearman--Brown reliability $r_K=2r_{\rm sh}/(1{+}r_{\rm sh})$ from the correlation
$r_{\rm sh}$ of two independent $K/2$-donor means. The \emph{information gain} $\mathcal G$ is the UAUC increase from adding
$\ell$ to the prior-only item mean $m(i)$ (other users' mean rating of $i$ strictly before the candidate's timestamp) in a
logistic stacker of $y$ on $(m,\ell)$, fitted on dev and evaluated on confirmation users. The \emph{popularity link}
$\mathcal L$ is the Spearman correlation of $\pi(i)$ with $\log(1{+}n_i)$ ($n_i$: all-time event count), also partialled on
$m(i)$.

\paragraph{Where can confidence change a ranking?}
\begin{lemma}[Invariance]\label{lem:invariance}\label{prop:monotone}
If $g_u:\mathbb R\to\mathbb R$ is strictly increasing and $s'_i=g_u(s_i)$ for $i\in C_u$, then
$\operatorname{sgn}(s'_i-s'_j)=\operatorname{sgn}(s_i-s_j)$ for all $i,j\in C_u$ (the weak order of $C_u$ is unchanged). Hence
$\mathrm{AUC}_u$, HR@$k$, NDCG@$k$ and MRR are unchanged: they depend on $s$ only through these signs (via $a$, $t$,
$c(s_i-s_j)$), under either tie convention (exact expectation over tie-breaking, or the metric at the expected rank
$1{+}a{+}t/2$).
\end{lemma}
\begin{proof}
Strict monotonicity preserves each of $s_i>s_j$, $s_i=s_j$ and $s_i<s_j$.
\end{proof}
Temperature scaling~\citep{guo2017calibration}, Platt scaling with positive slope, removal of $b_u$ and the list softmax
$w_i\propto\exp(\ell_i/T)$ are such maps: per-user calibration, however it is fitted, cannot change UAUC or within-unit
NDCG, let alone improve them. The lemma fails for merely non-decreasing maps (isotonic regression): they never reverse a
strict preference but create ties (a merged inverted pair gains $\frac12$, a merged correct pair loses $\frac12$), so a map
fitted on the evaluation labels can raise AUC through ties alone.

\begin{proposition}[Item-dependent corrections]\label{prop:itemdep}
For continuous, strictly increasing $g_{u,i}$ and $s'_i=g_{u,i}(s_i)$, $i\in C_u$, the weak order of $C_u$ is unchanged for
every score vector if and only if all $g_{u,i}$ coincide.
\end{proposition}
\begin{proof}
``If'' is Lemma~\ref{lem:invariance}. Conversely, let $g_{u,i}(x)<g_{u,j}(x)$ for some $i,j,x$. By continuity
$g_{u,i}(x{+}\varepsilon)<g_{u,j}(x)$ for some $\varepsilon>0$, so $s_j=x$, $s_i=x{+}\varepsilon$ gives $s_i>s_j$ but
$s'_i<s'_j$.
\end{proof}
Contextual~\citep{zhao2021calibrate} and batch calibration~\citep{zhou2024batch} subtract a prior estimated from
content-free inputs or from the batch; for a binary question this is a constant of the unit and rank-inert.
Candidate-conditioned priors (the no-history prompt, in the spirit of PMI scoring~\citep{holtzman2021surface}; the
history-swap prior $\hat\pi$) and second readouts of the same pair (a mirrored or paraphrased question) are
item-dependent: for $s'=s-\lambda b(i)$ a pair with $s_i>s_j$ is exchanged or tied exactly when
$s_i-s_j\le\lambda(b_i-b_j)$. Reordering is not improving, so each such correction is compared with an equal-compute
control. Across units the lemma is silent: per-user offsets change pooled AUC, thresholds and selective serving, and
training changes $f$ itself.

% PROP3-BEGIN
\begin{proposition}[Item prior and discriminability]\label{prop:contam}
Fix a unit; let $A(\cdot)$ be the population AUC, $\Pr(\text{a positive outscores a negative})$, of a score on its
candidates (UAUC estimates its mean over units). Assume that, given $y$, $\pi(i)$ and $e(u,i)$ are independent Gaussians
with class means $\nu_y,\mu_y$ (unit constants are absorbed in $b_u$) and class-independent standard deviations
$\tau_\pi,\tau_e$. With $d_e=(\mu_1-\mu_0)/\tau_e$, $d_\pi=(\nu_1-\nu_0)/\tau_\pi$ and $\rho=\tau_\pi^2/(\tau_\pi^2+\tau_e^2)$,
\[
  A(\ell)=\Phi\Big(\tfrac{1}{\sqrt2}\big[\sqrt{1-\rho}\,d_e+\sqrt{\rho}\,d_\pi\big]\Big),\qquad A(e)=\Phi\big(d_e/\sqrt2\big).
\]
If the prior carries no label information ($d_\pi=0$), then $\Phi^{-1}(A(\ell))=\sqrt{1-\rho}\,\Phi^{-1}(A(e))$.
\end{proposition}
\begin{proof}
For a positive $i$ and a negative $j$, $\ell(u,i)-\ell(u,j)$ ($b_u$ cancels) is Gaussian with mean
$\tau_\pi d_\pi+\tau_e d_e$ and variance $2(\tau_\pi^2+\tau_e^2)$; divide, using
$(\tau_\pi,\tau_e)/\sqrt{\tau_\pi^2+\tau_e^2}=(\sqrt\rho,\sqrt{1-\rho})$. Dropping $\pi$ gives $A(e)$.
\end{proof}
So a label-free item prior shrinks the sensitivity index of $e$ by $\sqrt{1-\rho}$, where $\rho$ is the item-prior share of
within-class variance (the analogue of $1-\mathcal S$). Subtracting a noisy $\hat\pi=\pi+\eta$ ($\eta$ Gaussian, independent
of $y$) leaves $e-\eta$, the case $d_\pi=0$ with $\tau_\eta$ for $\tau_\pi$, so the reliability of $\hat\pi$ must accompany any
prior-subtraction result. The binormal form is a benchmark, not a claim about the data.
% PROP3-END
